\color{unidarkgreen}{\subsection{Логика сигнализации}}
\color{black}

\begin{enumerate}

\item
Устройство формирует выходное воздействие на сигнализацию «Срабатывание» по схеме «ИЛИ» в следующих случаях:
\begin{itemize}
    \item срабатывание УРОВ В1(В) на себя;
    \item срабатывание УРОВ В1(В) на смежные выключатели.
\end{itemize}

\item
Устройство формирует выходное воздействие на сигнализацию «Неисправность» по схеме «ИЛИ» в следующих случаях:
\begin{itemize}
    \item неисправность ТТ;
    \item неисправность выключателя;
	\item неисправность ОТ терминала.
\end{itemize}

\item
Сброс выходных сигналов \textit{«Срабатывание фикс.»} и \textit{«Неисправность фикс.»}, выполненных с фиксацией производится подачей сигнала \textit{«Сброс»}.

\item
Функциональная схема блока логики сигнализации приведена на рисунке \ref{pict_sign:log_sign}.

\vspace{2mm}
\begin{figure}[H]
\centering
\includegraphics[width=0.8\textwidth,height=0.8\textheight,keepaspectratio]{\fbpath/10.02 АУВ/Сигнализация/img/log_sign.pdf}
\captionof{figure}{Функциональная схема блока логики сигнализации}\label{pict_sign:log_sign}
\end{figure}

\end{enumerate}